\VerdierPhysicalPageAnchor{0040}%
\VerdierSemanticAnchor{ch1-1-3-2-eq4}
et l'isomorphisme $c(g',g)^{-1}$ définit une application :

\begin{equation}
  \operatorname{Hom}_{\mathcal C}
    \bigl(T(g)Y,T(g)T(g')Z\bigr)
  \longrightarrow
  \operatorname{Hom}_{\mathcal C}
    \bigl(T(g)Y,T(g'g)Z\bigr).
  \tag{1.3.2.4}
\end{equation}

\VerdierSemanticAnchor{ch1-1-3-2-eq5}
D'où, en prenant l'image de $v$, un élément
$v'\in\operatorname{Hom}_{\mathcal C}(T(g)Y,T(g'g)Z)$. Le composé dans
$\mathcal C$ de $u$ et de $v'$ est un élément :

\begin{equation}
  w=v'\mathbin{\circ_{\mathcal C}}u,
  \qquad w\in\operatorname{Hom}^{g'g}(X,Z),
  \tag{1.3.2.5}
\end{equation}
qui est, par définition, le composé de $u$ et de $v$ dans
$\mathcal C[G]$. Le foncteur fibrant $\mathcal C[G]\to\mathcal G$ est
alors le foncteur évident. La catégorie fibrée
$\mathcal C[G]\to\mathcal G$ est munie d'un clivage normalisé canonique.

\VerdierSemanticAnchor{ch1-1-3-3}
\noindent\textbf{1.3.3.}\quad
On a en fait un résultat plus précis. Soient $\mathcal C$ et
$\mathcal C'$ deux catégories graduées de type $G$ munies de clivages
normalisés. Le foncteur $\Psi$ définit une bijection :

\VerdierSemanticAnchor{ch1-1-3-3-eq1}
\begin{equation}
  \Psi(\mathcal C,\mathcal C')\colon
  \operatorname{Hom}_{\operatorname{Grad}(G)}(\mathcal C,\mathcal C')
  \longrightarrow
  \operatorname{Hom}_{G^\circ\text{-}\operatorname{Cat}}
    \bigl(\Psi(\mathcal C),\Psi(\mathcal C')\bigr).
  \tag{1.3.3.1}
\end{equation}

Or ces ensembles sont les ensembles d'objets des catégories :

\VerdierSemanticAnchor{ch1-1-3-3-eq2}
\begin{equation}
  \mathcal{H}\!om_{\operatorname{Grad}(G)}(\mathcal C,\mathcal C'),
  \qquad
  \mathcal{H}\!om_{G^\circ\text{-}\operatorname{Cat}}
    \bigl(\Psi(\mathcal C),\Psi(\mathcal C')\bigr),
  \tag{1.3.3.2}
\end{equation}
dont les morphismes sont respectivement les morphismes de foncteurs de
degré $e$ et les morphismes de foncteurs compatibles avec les opérations
de $G^\circ$. De plus, il résulte immédiatement de (1.2.7) que
l'application $\Psi(\mathcal C,\mathcal C')$ s'étend en un foncteur :

\VerdierSemanticAnchor{ch1-1-3-3-eq3}
\begin{equation}
  \Psi(\mathcal C,\mathcal C')\colon
  \mathcal{H}\!om_{\operatorname{Grad}(G)}(\mathcal C,\mathcal C')
  \longrightarrow
  \mathcal{H}\!om_{G^\circ\text{-}\operatorname{Cat}}
    \bigl(\Psi(\mathcal C),\Psi(\mathcal C')\bigr)
  \tag{1.3.3.3}
\end{equation}
et on vérifie aisément :

\VerdierSemanticAnchor{ch1-1-3-4}
\noindent\textbf{Proposition 1.3.4.}\quad
\emph{Le foncteur $\Psi(\mathcal C,\mathcal C')$ est un isomorphisme de
catégories.}

\VerdierSemanticAnchor{ch1-1-3-5}
\noindent\textbf{1.3.5.}\quad
Enfin, pour être complet, il conviendrait de noter que le foncteur
$\Psi$ (étendu aux morphismes de foncteurs) respecte les convolutions des
foncteurs par les morphismes de foncteurs. La vérification en est laissée
au lecteur.

\VerdierSemanticAnchor{ch1-1-3-6}
\noindent\textbf{1.3.6.}\quad
Le plus souvent, nous utiliserons des $G^\circ$-catégories strictes
(1.2.2). Soient $(\mathcal C,G^\circ,T)$ et
$(\mathcal C',G^\circ,T')$ deux $G^\circ$-catégories strictes. Un
$G^\circ$-foncteur $(F,m)$ est alors un foncteur
$F\colon\mathcal C\to\mathcal C'$ muni d'isomorphismes :

\VerdierSemanticAnchor{ch1-1-3-6-eq1}
\begin{equation}
  m(g)\colon FT(g)\longrightarrow T'(g)F,
  \qquad g\in G,
  \qquad m(e)=\text{identité}
  \tag{1.3.6.1}
\end{equation}
